\section{Problem Formulation}

In this work, we formulate the sequential hyperparameter optimization problem as a Markov Decision Process (MDP). Let $\Lambda$ denote the $H$-dimensional hyperparameter search space. The MDP is defined by the tuple $\mathcal{M}\coloneqq (\mathcal{S}, \mathcal{A}, P, R, \gamma)$. The action space is defined as $\mathcal{A}\coloneqq\Lambda$, where an action $a_t\in\mathcal{A}$ corresponds to selecting a specific hyperparameter configuration $\lambda_t$ to evaluate. 

The state space $\mathcal{S}$ is defined as $\mathcal{S}\coloneqq\mathbb{R}^w\times(\Lambda\times R\times \mathcal{C})^*$, where each state $s_t\in\mathcal{S}$ encapsulates both the static dataset meta-features $\mathbf{d}\in\mathbb{R}^w$ and the dynamic history of previously evaluated configurations alongside their corresponding outcomes, including $\mathcal{C}\subset\mathbb{R}^K$ being the space of learning curves, representing a sequence of validation losses over $K$ training epochs and containing the information of the first and second derivatives of the learning curve. 

The reward function is defined as $R(D, a)=-f(D, \lambda=a)$ on a given dataset $D$. The deterministic transition function $P$ updates the state by appending the newly selected configuration and observed reward to the historical sequence, yielding $s_{t+1}=(s_t,(a_t, r_t, c_t))$. Furthermore, we impose strict termination constraints on the environment: an episode concludes either when a predefined budget sequence length $T$ is exhausted, or if the agent selects the identical action consecutively ($a_t=a_{t-1}$).

Our objective is to learn an optimal policy $\pi^*:\mathcal{S}\rightarrow\mathcal{A}$ that maximizes the expected discounted cumulative reward, effectively identifying the hyperparameter configuration that minimizes the generalization error across varying datasets. This is achieved by estimating the optimal action-value function $Q^*(s,a)$, which satisfies the Bellman equation:
$$Q^*(s, a)=\mathbb{E}_{\pi}\left[r+\gamma\max_{a'}Q^*(s',a')|S_0=s,A_0=a\right]$$
where $\gamma\in[0, 1]$ is the discount factor. To operate within this high-dimensional state space, which consists of variable-length temporal sequences and distinct response surfaces, we employ a parameterized function approximator $Q(s,a;\theta)$. Specifically, we utilize a Long Short-Term Memory (LSTM) network to model the dynamic action-value space, enabling the agent to capture long-range dependencies and transfer optimization knowledge across datasets.